%% Les six changements d'état, disposés sur un triangle « Gaz en haut ».
%%   - côté gauche  : Solide <-> Gaz      (sublimation / condensation)
%%   - côté droit   : Liquide <-> Gaz     (vaporisation / condensation-liquéfaction)
%%   - côté bas     : Solide <-> Liquide  (fusion / solidification)
%% Code couleur de l'unité : rouge (solideA) = le corps REÇOIT de l'énergie (Qp > 0),
%% bleu (solideB) = le corps CÈDE de l'énergie (Qp < 0).
%% Trois bulles montrent l'organisation microscopique des constituants.
\documentclass[tikz,border=4pt]{standalone}
\usepackage{liaisons}
\begin{document}
\begin{tikzpicture}[x=1cm,y=1cm,scale=0.85,
    etat/.style={draw=black!70, line width=0.9pt, rounded corners=3pt, fill=black!4,
                 inner sep=4pt, minimum width=1.5cm, font=\small\bfseries},
    fleche/.style={-{Stealth[length=5pt]}, line width=1.1pt},
    nom/.style={font=\scriptsize, inner sep=1pt, align=center},
    bulle/.style={draw=black!45, line width=0.6pt, fill=white},
    rappel/.style={draw=black!45, line width=0.4pt}]

%% ---------- Macro locale : une flèche parallèle à un côté, décalée de #4 --------------
%% #1 départ, #2 arrivée, #3 couleur, #4 décalage perpendiculaire (cm, signe = côté)
\newcommand\flcote[4]{%
  \pgfmathanglebetweenpoints{\pgfpointanchor{#1}{center}}{\pgfpointanchor{#2}{center}}%
  \edef\anglecote{\pgfmathresult}%
  \draw[fleche, #3] ([shift={(\anglecote+90:#4)}]$(#1)!0.19!(#2)$)
                 -- ([shift={(\anglecote+90:#4)}]$(#1)!0.81!(#2)$);}

%% ---------- Sommets du triangle ------------------------------------------------------
\node[etat] (G) at (0,3.0)     {Gaz};
\node[etat] (S) at (-3.1,0.15) {Solide};
\node[etat] (L) at (3.1,0.15)  {Liquide};

%% ---------- Côté gauche : solide <-> gaz ---------------------------------------------
\flcote{S}{G}{solideA}{0.30}                       % montante : sublimation
\flcote{G}{S}{solideB}{0.30}                       % descendante : condensation
\node[nom, text=solideA, rotate=42.6] at ([shift={(132.6:0.60)}]$(S)!0.5!(G)$) {Sublimation};
\node[nom, text=solideB, rotate=42.6] at ([shift={(-47.4:0.55)}]$(S)!0.55!(G)$) {Condensation};

%% ---------- Côté droit : liquide <-> gaz ---------------------------------------------
\flcote{L}{G}{solideA}{-0.30}                      % montante : vaporisation
\flcote{G}{L}{solideB}{-0.30}                      % descendante : liquéfaction
\node[nom, text=solideA, rotate=-42.6] at ([shift={(47.4:0.60)}]$(L)!0.5!(G)$) {Vaporisation};
\node[nom, text=solideB, rotate=-42.6] at ([shift={(-132.6:0.66)}]$(L)!0.60!(G)$)
  {Condensation\\(liquéfaction)};

%% ---------- Côté bas : solide <-> liquide --------------------------------------------
\flcote{S}{L}{solideA}{-0.30}                      % vers la droite, sous le côté : fusion
\flcote{L}{S}{solideB}{-0.30}                      % vers la gauche, au-dessus : solidification
\node[nom, text=solideA, anchor=north] at (0,-0.30) {Fusion};
\node[nom, text=solideB, fill=white, inner sep=2pt] at (0,0.45) {Solidification};

%% ---------- Bulles microscopiques ----------------------------------------------------
%% solide : réseau régulier (20 disques)
\draw[rappel] (S) -- (-2.55,-1.35);
\draw[bulle] (-2.55,-1.35) circle (0.62);
\begin{scope}[shift={(-2.55,-1.35)}]
  \foreach \i in {-2,-1,0,1,2} \foreach \j in {-1.5,-0.5,0.5,1.5}
    \fill[solideB!70!black] (0.21*\i,0.21*\j) circle (0.055);
\end{scope}
\node[font=\scriptsize, anchor=north] at (-2.55,-2.05) {réseau régulier};

%% liquide : disques jointifs mais désordonnés (8 disques)
\draw[rappel] (L) -- (2.55,-1.35);
\draw[bulle] (2.55,-1.35) circle (0.62);
\begin{scope}[shift={(2.55,-1.35)}]
  \foreach \p in {(-0.24,0.22),(0.05,0.30),(0.30,0.12),(-0.30,-0.08),(-0.02,0.00),
                  (0.25,-0.18),(-0.20,-0.30),(0.08,-0.32)}
    \fill[solideB!70!black] \p circle (0.10);
\end{scope}
\node[font=\scriptsize, anchor=north] at (2.55,-2.05) {jointifs, désordonnés};

%% gaz : 3 disques très espacés
\draw[rappel] (G) -- (2.05,3.35);
\draw[bulle] (2.05,3.35) circle (0.62);
\begin{scope}[shift={(2.05,3.35)}]
  \foreach \p in {(-0.28,0.24),(0.26,0.02),(-0.06,-0.32)}
    \fill[solideB!70!black] \p circle (0.085);
\end{scope}
\node[font=\scriptsize, anchor=west] at (2.72,3.35) {très espacés};

%% ---------- Légende ------------------------------------------------------------------
\node[font=\scriptsize, align=center, anchor=north] at (0,-2.55)
  {flèches \textcolor{solideA}{rouges} : le corps \textbf{reçoit} de l'énergie ($Q_p>0$)\\
   flèches \textcolor{solideB}{bleues} : il en \textbf{cède} ($Q_p<0$)};
\end{tikzpicture}
\end{document}
